\documentclass[9.5pt,a4paper]{article}
\usepackage[margin=0.5in,top=0.45in,bottom=0.45in]{geometry}
\usepackage{amsmath,amssymb}
\usepackage{booktabs}
\usepackage{tabularx}
\usepackage{microtype}
\usepackage{hyperref}
\usepackage{enumitem}

\hypersetup{
    colorlinks=true,
    linkcolor=blue,
    citecolor=blue,
    urlcolor=blue
}

\setlist{nosep,leftmargin=1.1em}

\title{\vspace{-1.2em}\textbf{\Large Beyond Token Serialization: Recurrent Latent Relaxation and In-Context Adaptation in BDH-CQ}\vspace{-0.7em}}
\author{}
\date{\vspace{-2.5em}}

\begin{document}

\maketitle
\thispagestyle{empty}

\begin{abstract}
\vspace{-0.4em}
\noindent Standard autoregressive Large Language Models (LLMs) rely on verbal Chain-of-Thought (CoT) prompting to execute multi-step reasoning. We argue that token serialization is an architectural crutch for the fixed computational depth of standard Transformers rather than a biological necessity. Pathway's Dragon Hatchling (BDH) architecture family, specifically BDH-CQ, provides a post-Transformer alternative: few-shot demonstrations are assimilated into recurrent associative state during the forward pass, and reasoning is performed via iterative latent relaxation. This briefing formalizes BDH-CQ mechanics, examines verified benchmark results on ARC-AGI-1 and Sudoku Extreme, and evaluates the non-trivial trade-off between compute efficiency and intermediate token-level interpretability.
\vspace{-0.4em}
\end{abstract}

\vspace{-0.6em}
\section{The Token Serialization Bottleneck}
\vspace{-0.3em}
Autoregressive Transformers have static layer depth $L$, strictly bounding the computation executed per forward pass. Chain-of-Thought (CoT) circumvents this constraint by externalizing intermediate steps as emitted natural-language tokens, feeding each symbol back into the attention window so subsequent passes borrow computational depth from sequence length.
However, token serialization introduces three fundamental structural inefficiencies:
\begin{enumerate}
    \item \textbf{Linear Memory Expansion:} State footprint scales as $\mathcal{O}(T)$ in token sequence length, inflating Key-Value (KV) cache memory.
    \item \textbf{Token Compute Overhead:} Full autoregressive projection is paid for conversational syntax, punctuation, and linguistic filler.
    \item \textbf{Verbosity-Bound Cost:} Inference FLOPs scale with natural language verbosity rather than intrinsic problem complexity.
\end{enumerate}

\vspace{-0.5em}
\section{Formal Mechanics of BDH-CQ}
\vspace{-0.3em}
BDH-CQ eliminates intermediate token generation by moving reasoning into continuous latent dynamics \cite{bdhcq2026, pathway2025}:

\vspace{0.15em}
\noindent\textbf{1. Forward-Pass Demonstration Assimilation.} Unlike test-time optimization frameworks that execute gradient descent across demonstration samples, BDH-CQ updates an associative context state $S_t$ strictly during the inference forward pass:
\begin{equation}
    S_t = U_{\theta}(S_{t-1}, D_t)
\end{equation}
where $D_t$ is an observed input-output demonstration and $U_\theta$ is a learned, frozen update operator. Context is compressed into fixed-rank associative memory using linear attention and ReLU low-rank projections, maintaining $\mathcal{O}(1)$ working memory complexity \cite{bdhcq2026}.

\vspace{0.15em}
\noindent\textbf{2. Recurrent Latent Relaxation.} Queries are resolved by repeatedly cycling a continuous hidden vector $h \in \mathbb{R}^d$ through a recurrent operator across depth $t \in \{1, \dots, T\}$ without token generation:
\begin{equation}
    h_{t+1} = \mathrm{LayerNorm}\big(h_t + \alpha \cdot \mathrm{ReLU}(W h_t - b)\big)
\end{equation}
The weight matrix $W$ and bias $b$ define a dynamical system whose attractor basins correspond to constraint satisfaction, constrained by biologically inspired non-negative activation sparsity ($\sim$5\% active units) \cite{pathway2025}. Sequence length remains invariant ($\mathcal{O}(1)$).

\vspace{-0.4em}
\section{Architectural Landscape Comparison}
\vspace{-0.3em}
Table~\ref{tab:comparison} contrasts BDH-CQ against token-based CoT and optimization solvers on the ARC-AGI benchmark \cite{chollet2019}.

\begin{table}[h!]
\centering
\footnotesize
\setlength{\tabcolsep}{3.5pt}
\renewcommand{\arraystretch}{1.05}
\begin{tabularx}{\textwidth}{@{}lXXX@{}}
\toprule
\textbf{Dimension} & \textbf{Standard Dense Transformer + CoT} & \textbf{Test-Time Optimization (HRM / TRM)} & \textbf{BDH-CQ (Pathway Research)} \\ \midrule
\textbf{Sequence Length} & $\mathcal{O}(T)$ (expands per reasoning token) & $\mathcal{O}(1)$ (fixed prompt structure) & $\mathcal{O}(1)$ (invariant continuous state) \\
\textbf{KV-Cache Scaling} & $\mathcal{O}(T)$ (linearly growing memory) & $\mathcal{O}(1)$ (no token accumulation) & $\mathcal{O}(1)$ (flat constant buffer) \\
\textbf{Test-Time Adaptation} & Appended prompt context tokens & Test-time backpropagation / weight updates & Forward-pass associative state updates ($U_\theta$) \\
\textbf{Compute Mechanism} & Autoregressive sampling per step & Multi-epoch backward-pass search & Low-rank recurrent forward relaxation \\
\textbf{Observability} & High (human-readable English tokens) & Low (parameter distribution shift) & Low (continuous vector trajectory in $\mathbb{R}^d$) \\ \bottomrule
\end{tabularx}
\caption{Architectural comparison of test-time reasoning paradigms.}
\label{tab:comparison}
\end{table}

\vspace{-1.1em}
\section{Empirical Evidence \& Pareto Dynamics}
\vspace{-0.3em}
Verified benchmarks establish that recurrent latent computation shifts the ``intelligence-per-dollar'' Pareto frontier:
\begin{itemize}
    \item \textbf{ARC-AGI-1 Benchmark \cite{chollet2019}:} A 150M parameter BDH-CQ model achieved \textbf{29.5\% pass@2} at a verified inference cost of \textbf{\$0.0007 per task} \cite{bdhcq2026}. By comparison, GPT-5.6 Luna Low scored 34.2\% at \$0.0077 per task---an 11$\times$ cost premium for a 4.7 percentage point margin. Independent evaluations (including replications by \L{}. Kaiser and R. Zhong) verify that forward-pass associative memory reliably acquires inductive rules without test-time weight updates.
    \item \textbf{Sudoku Extreme Evaluation:} Over $\sim$250,000 extreme puzzles, Pathway BDH attained \textbf{97.4\% accuracy} without emitting verbal tokens, executing backtracking routines, or relying on external symbolic verifiers \cite{pathway2025}, where dense LLMs (o3-mini, DeepSeek R1) collapsed due to cumulative token drift.
\end{itemize}

\vspace{-0.5em}
\section{Scientific Caveats \& Honest Limitations}
\vspace{-0.3em}
Despite compute efficiency gains, recurrent latent relaxation incurs trade-offs that bound its current utility:
\begin{enumerate}
    \item \textbf{The Observability Gap:} CoT generates an auditable text scratchpad. When BDH-CQ settles into an incorrect attractor, no token-level trace exists to diagnose whether failure stemmed from demonstration encoding or recurrent relaxation dynamics.
    \item \textbf{Non-Convex Attractor Trapping:} Recurrent relaxation is an energy minimization process. On complex coupled constraints, $h_t$ can stall in shallow local minima without the explicit backtracking pathways available to symbolic or verbalized solvers.
    \item \textbf{Domain Bounds:} Verified advantages are proven on discrete constraint satisfaction (ARC-AGI, Sudoku); generalization to open-ended dialogue and external tool use remains unestablished.
\end{enumerate}

\vspace{-0.6em}
\begin{thebibliography}{3}
\footnotesize
\setlength{\itemsep}{-0.5pt}
\bibitem{bdhcq2026} B. Engdahl, J. Chorowski, A. Kosowski, Z. Stamirowska, et al. \textit{BDH-CQ: In-Context Learning with Recurrent Latent Reasoning}. arXiv:2608.09888, 2026.
\bibitem{pathway2025} Pathway Research. \textit{BDH (Dragon Hatchling): A Brain-Inspired Post-Transformer Architecture}. Technical Report, 2025/2026.
\bibitem{chollet2019} F. Chollet. \textit{On the Measure of Intelligence}. arXiv:1911.01547, 2019.
\end{thebibliography}

\end{document}